\documentclass[10pt,a4paper]{amsart}
\usepackage[margin=19mm]{geometry}
\usepackage{amsmath,amssymb,mathtools}
\usepackage[scr=rsfs,cal=euler]{mathalfa}
\usepackage{xfrac}
\usepackage[unicode,hidelinks,bookmarks=false]{hyperref}
\newcommand{\ie}{i.e.\ }
\newcommand{\cf}{cf.\ }
\newcommand{\resp}{respectively\ }
\newcommand{\Subsection}{\subsection}
\providecommand{\nobreakdash}{\nobreak\hbox{-}}
\setlength{\emergencystretch}{2em}
\multlinegap0pt
\usepackage{amssymb}
\usepackage{cancel}
\usepackage{tikz}
\usepackage[shortlabels]{enumitem}
\usepackage{float}
\usepackage{bm}
\newtheorem{lemma}{Lemma}
\newcommand{\R}{\mathbb{R}}
\newcommand{\abs}[1]{\lvert#1\rvert}
\title{A partial answer to Brezis' Open Problem 2.1}
\author{Chao Ji, Kai Sheng}
\date{Source: arXiv:2608.21896v1 (CC BY 4.0)}
\begin{document}
\maketitle

\noindent\emph{Source and licence.} A partial answer to Brezis' Open Problem 2.1. Version arXiv:2608.21896v1 (CC BY 4.0), distributed by arXiv under the Creative Commons Attribution 4.0 International licence, http://creativecommons.org/licenses/by/4.0/, as recorded in the arXiv licence metadata for this exact version and shown on the abstract page https://arxiv.org/abs/2608.21896v1. Full licence notice: \texttt{licenses/arxiv-2608.21896-CC-BY-4.0.txt}.\par\smallskip

\noindent\textit{Editorial scope.} Counted unit: the article's lemma lem:convexity-threshold (source0.tex lines 541--546). The article's complete original proof of it is delivered with it (source0.tex lines 548--665, one \texttt{proof} environment of 118 lines). No mathematics is written, repaired or re-derived: the statement and the proof are the article's own lines, in the article's own notation, and no retained line is substituted or re-flowed.  Retained uncounted as the source-local context the proof needs: lines 227--237; lines 363--380; lines 411--415.  Nothing was omitted from any retained span, so the package carries no omission record.  No retained line contains a citation, so the wrapper carries no bibliography and no bibliography entry.  No retained line contains a figure environment, an image-inclusion command or drawing code, so no image is delivered and the package declares no asset dependency.  Editorial formatting only, in addition: an AMS \texttt{amsart} preamble that restates the article's own declarations and macros used by the retained lines, namely the environments \texttt{lemma} on separate counters, together with the standard packages those lines need.  The retained environments print in this document as Lemma 1 in order of appearance, whereas the article prints the same items with the numbering of its own full document.  The complete native source of the article as distributed in the arXiv e-print for this exact version is bundled unchanged as \texttt{sources/208284/source0.tex}.
\begin{equation}\label{eq:general-GL-system}
	\begin{cases}
		-\Delta u
		=
		\dfrac{1}{\varepsilon^2}
		u\bigl(1-\abs{u}^2\bigr)
		& \text{in } \Omega,\\[2mm]
		u=g
		& \text{on } \partial\Omega,
	\end{cases}
\end{equation}

\begin{equation}\label{eq:second-variation}
	\begin{aligned}
		D^2E_\varepsilon(u)[\varphi,\psi]
		={}&
		\int_\Omega
		\nabla\varphi:\nabla\psi
		-
		\frac{1}{\varepsilon^2}
		\int_\Omega
		\bigl(1-\abs{u}^2\bigr)\varphi\cdot\psi\\
		&+
		\frac{2}{\varepsilon^2}
		\int_\Omega
		(u\cdot\varphi)(u\cdot\psi)
	\end{aligned}
	\qquad
	\forall\,\varphi,\psi\in H_0^1(\Omega;\R^m).
\end{equation}

\begin{lemma}\label{lem:existence}
	For every $\varepsilon>0$, problem
	\eqref{eq:general-GL-system} admits a weak solution
	$u\in H_g^1(\Omega;\R^m)$.
\end{lemma}

\begin{lemma}\label{lem:convexity-threshold}
	For every $\varepsilon\geq\varepsilon_*$, the functional
	$E_\varepsilon$ is strictly convex on
	$H_g^1(\Omega;\R^m)$. Consequently,
	problem \eqref{eq:general-GL-system} has a unique weak solution.
\end{lemma}

\begin{proof}
	Let $u\in H_g^1(\Omega;\R^m)$ and
	$h\in H_0^1(\Omega;\R^m)\setminus\{0\}$. By
	\eqref{eq:second-variation},
	\[
		\begin{aligned}
		D^2E_\varepsilon(u)[h,h]
		={}&
		\int_\Omega\abs{\nabla h}^2
		-
		\frac{1}{\varepsilon^2}
		\int_\Omega\abs{h}^2\\
		&+
		\frac{1}{\varepsilon^2}
		\int_\Omega\abs{u}^2\abs{h}^2
		+
		\frac{2}{\varepsilon^2}
		\int_\Omega(u\cdot h)^2.
		\end{aligned}
	\]
	Since $\varepsilon\geq\varepsilon_*$, the Poincar\'e inequality gives
	\begin{equation}\label{eq:convexity-poincare}
		\int_\Omega\abs{\nabla h}^2
		-
		\frac{1}{\varepsilon^2}
		\int_\Omega\abs{h}^2
		\geq
		\left(
		\lambda_1-\frac{1}{\varepsilon^2}
		\right)
		\int_\Omega\abs{h}^2
		\geq0.
	\end{equation}
	Suppose that $D^2E_\varepsilon(u)[h,h]=0$. By
	\eqref{eq:second-variation} and \eqref{eq:convexity-poincare},
	\[
		\int_\Omega\abs{\nabla h}^2
		-
		\frac{1}{\varepsilon^2}
		\int_\Omega\abs{h}^2
		=
		0,
		\qquad
		\int_\Omega\abs{u}^2\abs{h}^2
		=
		0.
	\]
	Since $h\neq0$, \eqref{eq:convexity-poincare} gives
	$\varepsilon=\varepsilon_*$ and
	\[
		\int_\Omega\abs{\nabla h}^2
		=
		\lambda_1\int_\Omega\abs{h}^2.
	\]
	Since $\Omega$ is connected, the first Dirichlet eigenvalue
	$\lambda_1$ is simple. Moreover, for every
	$w=(w_1,\ldots,w_m)\in H_0^1(\Omega;\R^m)$,
	\[
		\int_\Omega\abs{\nabla w}^2
		-
		\lambda_1\int_\Omega\abs{w}^2
		=
		\sum_{\alpha=1}^m
		\left(
		\int_\Omega\abs{\nabla w_\alpha}^2
		-
		\lambda_1\int_\Omega\abs{w_\alpha}^2
		\right).
	\]
	Each term on the right-hand side is nonnegative by the Poincar\'e
	inequality. Therefore,
	\begin{equation}\label{eq:first-eigenspace-characterization}
		\int_\Omega\abs{\nabla w}^2
		=
		\lambda_1\int_\Omega\abs{w}^2
		\quad\Longrightarrow\quad
		w=\phi_1\xi
		\quad\text{for some }\xi\in\R^m,
	\end{equation}
	where $\phi_1$ is the positive Dirichlet eigenfunction corresponding
	to $\lambda_1$. Applying
	\eqref{eq:first-eigenspace-characterization} to $h$, we have
	$h=\phi_1\xi$ for some $\xi\in\R^m\setminus\{0\}$. It follows that
	\[
		0
		=
		\int_\Omega\abs{u}^2\abs{h}^2
		=
		\abs{\xi}^2
		\int_\Omega\abs{u}^2\phi_1^2.
	\]
	Since $\phi_1>0$ in $\Omega$, we obtain $u=0$ a.e. in $\Omega$,
	which contradicts $u\in H_g^1(\Omega;\R^m)$ and $g\not\equiv0$.
	Therefore $D^2E_\varepsilon(u)[h,h]>0$ for every
	$u\in H_g^1(\Omega;\R^m)$ and every
	$h\in H_0^1(\Omega;\R^m)\setminus\{0\}$. Let
	$u_0,u_1\in H_g^1(\Omega;\R^m)$ with $u_0\neq u_1$, and define
	$q\colon[0,1]\to\R$ by
	$q(t):=E_\varepsilon(u_0+t(u_1-u_0))$. Then
	\[
		q''(t)
		=
		D^2E_\varepsilon
		\bigl(u_0+t(u_1-u_0)\bigr)
		[u_1-u_0,u_1-u_0]
		>
		0
	\]
	for every $t\in[0,1]$. Hence $q$ is strictly convex, and therefore
	$E_\varepsilon$ is strictly convex on $H_g^1(\Omega;\R^m)$.
	By Lemma~\ref{lem:existence}, a weak solution exists for every
	$\varepsilon>0$. Moreover, weak solutions are precisely the critical
	points of $E_\varepsilon$ on $H_g^1(\Omega;\R^m)$, and a
	differentiable strictly convex functional has at most one critical
	point. Consequently, \eqref{eq:general-GL-system} has a unique weak
	solution for every $\varepsilon\geq\varepsilon_*$.

\end{proof}

\end{document}
